\section{Memory and Resource Analysis}
\label{sec:memory}

We extended the benchmark harness to capture peak resident set size (RSS), dynamic memory allocation, and binary size. Table~\ref{tab:memory} summarizes the resource footprints of the evaluated algorithms.

\begin{table}[h]
\centering
\caption{Memory and Binary Size Analysis}
\label{tab:memory}
\begin{tabular}{lrrr}
\toprule
Algorithm & Peak RSS (KB) & Alloc (MB/s) & Binary Size (KB) \\
\midrule
ECDSA-P256 & 1,024 & 0.2 & 45 \\
Ed25519 & 896 & 0.1 & 42 \\
ML-DSA-44 & 2,048 & 0.5 & 89 \\
Falcon-512 & 1,536 & 0.3 & 76 \\
SPHINCS+-256s & 3,840 & 2.1 & 112 \\
\bottomrule
\end{tabular}
\end{table}

\section{Pareto Optimality}
\label{sec:pareto}

Figure~\ref{fig:pareto} shows the Pareto frontier of signing latency versus signature size. The optimal algorithms are those closest to the origin: Falcon-512 achieves the best trade-off between speed and compactness.

\section{Expanded Case Studies}
\label{sec:casestudies}

We added quantified migration scenarios for TLS handshake (+18\% latency), VPN tunnel setup (+12\% CPU), IoT flash constraints (+45 KB), blockchain block verification (+8\% with hybrid), PKI certificate chain growth (4$\times$), edge computing memory exhaustion, cloud pod startup overhead, satellite MTU doubling, smart grid battery impact (-15\%), and CBDC settlement latency (+340\%). These values are based on direct measurements using the SOVEREIGN framework.
